% Nová příprava: 2026-10-28; 1. ročník SPŠ, etapa 13–24.
% Rozpis: 1_rocnik/tematicke_plany/Rozpis_vyuky_MAT_1_rocnik.tex, hodina 24.
% Původ: vlastní formulace výkladu a vlastní úlohy; metodická inspirace uvedenými zdroji.
% Kontrola, obsahové porovnání zdrojů a časy: planovani/DOKONCENI_PRIPRAV_13_24.md.
% Učebnice: ucebnice/03_Planimetrie/Planimetrie.tex, Trojúhelník a značení;
% ucebnice/04_Goniometrie_a_trigonometrie/Goniometrické funkce.tex,
% Tabulkové hodnoty (pouze ostré úhly) a praktické použití inverzních funkcí.
% Realisticky.cz: externi_zdroje/realisticky/01_Zakladni_poznatky/01_Opakovani/010119_Podobnosti_trojuhelniku_goniometricke_funkce/010119_Lekce.pdf
% Realisticky.cz: externi_zdroje/realisticky/01_Zakladni_poznatky/01_Opakovani/010120_Pouziti_goniometrickych_funkci_ve_fyzice/010120_Lekce.pdf
\documentclass[11pt,a4paper]{article}
% Spolecny styl pro ucitelske pripravy MAT (Work/MAT)
% Verze 1.1, 2026-10-23
\NeedsTeXFormat{LaTeX2e}
\ProvidesPackage{priprava}[2026/10/23 v1.1 Sablona ucitelskych priprav MAT]

% Zakladni typografie a cestina
\RequirePackage[T1]{fontenc}
\RequirePackage[utf8]{inputenc}
\RequirePackage[czech]{babel}
\RequirePackage{lmodern}
\RequirePackage{microtype}
\RequirePackage{amsmath,amssymb}
\RequirePackage{xcolor}
\RequirePackage{enumitem}
\RequirePackage[most]{tcolorbox}
\RequirePackage{tabularx,booktabs,array}
\RequirePackage{fancyhdr}
\RequirePackage{needspace}
\RequirePackage{geometry}
\RequirePackage[hidelinks]{hyperref}

\geometry{a4paper,top=17mm,bottom=18mm,left=23mm,right=23mm}
\definecolor{notegray}{gray}{0.39}
\setlength{\parindent}{0pt}
\setlength{\parskip}{2.5pt plus 1pt minus 0.5pt}
\setlist[itemize]{topsep=2pt,itemsep=1pt,parsep=0pt,leftmargin=1.45em}
\setlist[enumerate]{topsep=2pt,itemsep=2pt,parsep=0pt,leftmargin=1.65em}
\setcounter{secnumdepth}{0}
\renewcommand{\arraystretch}{1.12}

% Zahlavi: vlevo rocnik | globalni tema, vpravo cislo: nazev hodiny
% Zapati zustava zcela prazdne.
\pagestyle{fancy}
\fancyhf{}
\renewcommand{\headrulewidth}{0.35pt}
\renewcommand{\footrulewidth}{0pt}
\setlength{\headheight}{14pt}
\newcommand{\nastavPripravu}[4]{%
  \fancyhead[L]{\small #1\enspace\textbar\enspace #2}%
  \fancyhead[R]{\small #3: #4}%
}

% Minimalisticke cernobile boxy. Kazda definice zvlast.
\tcbset{commonlesson/.style={enhanced,breakable=false,colback=white,colframe=black,
  colbacktitle=white,coltitle=black,boxrule=0.45pt,
  arc=0.8mm,outer arc=0.8mm,left=3mm,right=3mm,top=1.6mm,bottom=1.6mm,
  toptitle=1mm,bottomtitle=0.7mm,fonttitle=\bfseries,
  attach boxed title to top left={xshift=2.5mm,yshift=-1.15mm},
  boxed title style={colback=white,colframe=white,boxrule=0pt,left=1mm,right=1mm}}}
\newtcolorbox{zadaniBox}[1][]{commonlesson,title={Zadání},#1}
\newtcolorbox{prikladBox}[1][]{commonlesson,title={Příklad},#1}
\newtcolorbox{definiceBox}[1][]{commonlesson,title={Definice},#1}

% Spolecny box pro zadani a jeho reseni, bez mezery mezi castmi.
% Pouziti:
%   \begin{prikladReseniBox}
%     Zadani ulohy.
%     \tcblower
%     \begin{align*} ... \end{align*}
%   \end{prikladReseniBox}
% Spodni cast ma jemny sedy podklad a delici linku; lze v ni sazet
% vice radku matematiky. Obe casti maji jen jeden spolecny ramecek.
\tcbset{solutionlesson/.style={
  commonlesson,
  enhanced jigsaw,
  breakable,
  colbacklower=black!3,
  segmentation style={solid,draw=black!25,line width=0.35pt},
  fontlower=\small\color{notegray},
  before lower={\noindent\textbf{Řešení:}\enspace},
  leftlower=3mm,rightlower=3mm
}}
\newtcolorbox{prikladReseniBox}[1][]{solutionlesson,title={Příklad},#1}
\newtcolorbox{zadaniReseniBox}[1][]{solutionlesson,title={Zadání},#1}

% Pomucky pro autora pripravy: sedive poznamky; \reseni pro starsi pripravy
\newcommand{\poznamka}[1]{\par\begingroup\small\color{notegray}\itshape
  \textnormal{\textbf{Poznámka pro učitele.}}\enspace #1\par\endgroup}
\newcommand{\reseni}[1]{\par\begingroup\small\color{notegray}
  \textbf{Očekávané řešení:}\enspace #1\par\endgroup}
\newcommand{\blok}[1]{\needspace{5\baselineskip}\subsection*{#1}}
\newcommand{\vyrok}[1]{\uv{#1}}

\usepackage{tikz}
\usetikzlibrary{patterns,calc,angles,quotes}
\nastavPripravu{1. ročník}{Goniometrie}{24}{Výšky a vzdálenosti}
\begin{document}
\blok{Ze situace udělat pravoúhlý trojúhelník}
Nejprve rozhodni, která délka je vodorovná, která svislá a která šikmá.
Zakresli pravý úhel a úhel zadaný; potom teprve vyber sinus, kosinus nebo tangens.
Používáme ideální model přímých úseček, vodorovného terénu a svislých objektů,
pokud zadání neříká jinak.
\begin{definiceBox}
\textbf{Výškový úhel} je úhel mezi vodorovným směrem a záměrnou přímkou
směřující vzhůru k pozorovanému bodu.
\end{definiceBox}
\begin{definiceBox}
\textbf{Hloubkový úhel} je úhel mezi vodorovným směrem a záměrnou přímkou
směřující dolů. Počítáme jeho kladnou velikost.
\end{definiceBox}
Kalkulačku nastav na stupně. Délky převáděj na stejné jednotky před dosazením,
mezivýsledky ponech nezaokrouhlené. Výsledné délky zaokrouhluj na setiny metru,
úhly na desetiny stupně; desetinná místa sama nezaručují přesnost reálného měření.

\begin{prikladReseniBox}[title={P1: výška stavby a výška očí}]
Pozorovatel stojí na vodorovném terénu $24\,\mathrm m$ od paty svislé stavby.
Oči má $1{,}60\,\mathrm m$ nad úrovní její paty. Vrchol vidí pod výškovým úhlem
$32^\circ$. Urči výšku stavby $H$.
\tcblower
\begin{center}
\begin{tikzpicture}[scale=.75,color=black]
\draw[thick] (-.8,0)--(5.2,0);
\draw[thick] (4,0)--(4,2.7667);
\draw (0,0)--(0,.2667); \fill (0,.2667) circle(2pt);
\draw[dashed] (0,.2667)--(4,.2667);
\draw[<->] (0,-.55)--node[below]{$d=24\,\mathrm m$}(4,-.55);
\draw[thick] (0,.2667)--(4,2.7667);
\draw (3.75,.2667)--(3.75,.5167)--(4,.5167);
\draw (.7,.2667) arc(0:32:.7);
\node at (1.05,.6267) {$32^\circ$};
\draw[<->] (-.5,0)--node[left]{$1{,}60\,\mathrm m$}(-.5,.2667);
\draw[<->] (4.65,0)--node[right]{$H$}(4.65,2.7667);
\end{tikzpicture}
\end{center}
V pravoúhlém trojúhelníku je svislou odvěsnou rozdíl $H-1{,}60$, nikoli celé $H$.
\[
\tan32^\circ=\frac{H-1{,}60}{24},
\qquad H=1{,}60+24\tan32^\circ\approx16{,}60\,\mathrm m.
\]
Kontrola: bez výšky očí by vyšla jen výška nad vodorovnou záměrou.
Úhel je menší než $45^\circ$, proto je tento výškový rozdíl menší než $24\,\mathrm m$.
\end{prikladReseniBox}

\begin{prikladReseniBox}[title={P2: šikmá délka není vodorovná vzdálenost},breakable=false]
Přímé napnuté lano délky $20\,\mathrm m$ svírá s vodorovnou rovinou úhel $65^\circ$.
Jeho dolní konec je na zemi. Urči výšku horního konce a vodorovnou vzdálenost konců.
\tcblower
Lano je přepona $L=20\,\mathrm m$. Výška $h$ je proti zadanému úhlu,
vodorovná vzdálenost $d$ je přilehlá:
\[
h=L\sin65^\circ\approx18{,}13\,\mathrm m,\qquad
d=L\cos65^\circ\approx8{,}45\,\mathrm m.
\]
Obě odvěsny jsou kratší než lano. Z nezaokrouhlených hodnot ověříme
$h^2+d^2=20^2$. Výraz $20\tan65^\circ$ by chybně považoval lano za odvěsnu.
\end{prikladReseniBox}

\blok{Navazující příklad podle času}
\begin{prikladReseniBox}[title={P3: potřebujeme dosah i délku rampy},breakable=false]
Rampa má překonat výškový rozdíl $900\,\mathrm{mm}$ pod úhlem $7^\circ$
k vodorovné rovině. Urči vodorovný dosah $d$ a délku rampy $L$.
\tcblower
Nejprve $900\,\mathrm{mm}=0{,}9\,\mathrm m$.
\[
\tan7^\circ=\frac{0{,}9}{d}\Rightarrow
d=\frac{0{,}9}{\tan7^\circ}\approx7{,}33\,\mathrm m,
\]
\[
\sin7^\circ=\frac{0{,}9}{L}\Rightarrow
L=\frac{0{,}9}{\sin7^\circ}\approx7{,}38\,\mathrm m.
\]
Šikmá délka $L$ musí být větší než vodorovný dosah $d$.
Oba výsledky porovnej s náčrtem: přepona není vodorovný dosah.
\end{prikladReseniBox}
\poznamka{Jádro P1, P2 a S1 a), b), přibližně 38 min práce,
organizace a závěr 7 min. P3 a S1 c), d) jsou další nabídka k výběru
nebo domů. Před každým výpočtem žák nakreslí vlastní trojúhelník.
Ve S1 b) zdůvodni shodu hloubkového a příslušného výškového úhlu
pomocí rovnoběžných vodorovných přímek.}

\blok{Samostatné procvičení}
\begin{zadaniReseniBox}[title={S1: model, vzorec, jednotka, kontrola},breakable=false]
\begin{enumerate}[label=\alph*)]
\item Z místa $18\,\mathrm m$ vodorovně od paty věže vidíme vrchol
pod výškovým úhlem $40^\circ$. Oči jsou $1{,}50\,\mathrm m$ nad úrovní paty. Urči výšku.
\item Z bodu $12\,\mathrm m$ nad vodorovnou hladinou vidíme loď pod hloubkovým úhlem
$38^\circ$. Urči vodorovnou vzdálenost lodě od svislého průmětu pozorovatele.
\item Přímá vzpěra délky $350\,\mathrm{cm}$ svírá s vodorovnou rovinou úhel $52^\circ$.
Urči svislý a vodorovný rozdíl poloh jejích konců v metrech.
\item Jak se liší výška při úhlu $30^\circ$, je-li $50\,\mathrm m$
vodorovnou vzdáleností, a je-li $50\,\mathrm m$ šikmou délkou?
Dolní záměrný bod má v obou modelech stejnou výšku jako základna.
\end{enumerate}
\tcblower
a) $H=1{,}50+18\tan40^\circ\approx16{,}60\,\mathrm m$.
b) Odpovídající ostrý úhel u hladiny má také $38^\circ$.
$d=\frac{12}{\tan38^\circ}\approx15{,}36\,\mathrm m$.
c) $350\,\mathrm{cm}=3{,}5\,\mathrm m$;
$h=3{,}5\sin52^\circ\approx2{,}76\,\mathrm m$,
$d=3{,}5\cos52^\circ\approx2{,}15\,\mathrm m$.
d) První model: $h=50\tan30^\circ\approx28{,}87\,\mathrm m$.
Druhý: $h=50\sin30^\circ=25\,\mathrm m$.
Označení délky nelze nahradit samotným číslem.
\end{zadaniReseniBox}

\blok{Volitelná rezerva}
\begin{zadaniReseniBox}[title={R1: sklon v procentech, 8–12 minut},breakable=false]
V této úloze sklon $p$ procent znamená $p=100h/d$,
kde $h$ je převýšení a $d$ vodorovný dosah.
Přímý úsek má sklon $8\,\%$ a vodorovný dosah $12\,\mathrm m$.
Urči převýšení, úhel sklonu a skutečnou šikmou délku.
\tcblower
$h=0{,}08\cdot12=0{,}96\,\mathrm m$.
$\tan\alpha=h/d=0{,}08$, tedy $\alpha=\arctan0{,}08\approx4{,}6^\circ$.
$L=\sqrt{12^2+0{,}96^2}\approx12{,}04\,\mathrm m$.
Sklon $8\,\%$ tedy neznamená úhel $8^\circ$.
\end{zadaniReseniBox}
\begin{zadaniReseniBox}[title={R2: citlivost měření, 5–8 minut},breakable=false]
V P1 se úhel může pohybovat od $31^\circ$ do $33^\circ$.
Jaké krajní výšky dostaneme při jinak stejných údajích?
\tcblower
$H_{\min}=1{,}60+24\tan31^\circ\approx16{,}02\,\mathrm m$,
$H_{\max}=1{,}60+24\tan33^\circ\approx17{,}19\,\mathrm m$.
Při pevném vodorovném dosahu většímu ostrému úhlu odpovídá větší výškový rozdíl.
Zápis výsledku na centimetry proto neznamená, že skutečnou výšku známe na centimetr.
\end{zadaniReseniBox}

\blok{Domácí procvičení}
\begin{zadaniReseniBox}[title={D1: 15–20 minut; každý model doplň náčrtem}]
\begin{enumerate}[label=\alph*)]
\item Stavbu pozorujeme z vodorovné vzdálenosti $30\,\mathrm m$ pod úhlem $35^\circ$.
Oči jsou $1{,}70\,\mathrm m$ nad úrovní její paty. Urči výšku stavby.
\item Lano délky $12\,\mathrm m$ svírá s vodorovnou rovinou $70^\circ$.
Urči svislý a vodorovný rozdíl poloh konců.
\item Přímý úsek má sklon $10\,\%$ ve smyslu $100h/d$
a vodorovný dosah $7{,}5\,\mathrm m$. Urči převýšení a úhel.
\item V rovině jsou body $A,B$ na přímém břehu řeky, $AB=12\,\mathrm m$.
Strom $C$ na protějším břehu leží na kolmici k $AB$ vedené bodem $B$.
Je-li $\angle CAB=56^\circ$, urči šířku $BC$.
\end{enumerate}
\tcblower
a) $H=1{,}70+30\tan35^\circ\approx22{,}71\,\mathrm m$.
b) $h=12\sin70^\circ\approx11{,}28\,\mathrm m$,
$d=12\cos70^\circ\approx4{,}10\,\mathrm m$.
c) $h=0{,}10\cdot7{,}5=0{,}75\,\mathrm m$,
$\alpha=\arctan0{,}10\approx5{,}7^\circ$.
d) Pravý úhel je u $B$; $AB$ je přilehlá odvěsna,
$BC$ protilehlá. Proto $BC=12\tan56^\circ\approx17{,}79\,\mathrm m$.
\end{zadaniReseniBox}
\poznamka{Domácí c) užívá definici sklonu přímo ze zadání, takže nevyžaduje
probrání volitelné rezervy. Navazující hodina 25 zpřesní význam mocnin;
dosavadní druhé mocniny a odmocniny byly konkrétní prostředky geometrických výpočtů.}
\end{document}
